\documentclass[10pt,a4paper]{amsart}
\usepackage[margin=19mm]{geometry}
\usepackage{amsmath,amssymb,mathtools}
\usepackage[scr=rsfs,cal=euler]{mathalfa}
\usepackage{xfrac}
\usepackage[unicode,hidelinks,bookmarks=false]{hyperref}
\newcommand{\ie}{i.e.\ }
\newcommand{\cf}{cf.\ }
\newcommand{\resp}{respectively\ }
\newcommand{\Subsection}{\subsection}
\providecommand{\nobreakdash}{\nobreak\hbox{-}}
\setlength{\emergencystretch}{2em}
\multlinegap0pt
\usepackage{amssymb}
\usepackage{mathtools}
\usepackage{xcolor}
\usepackage[shortlabels]{enumitem}
\usepackage{tikz}
\usepackage{csquotes}
\newtheorem{lemma}{Lemma}
\newtheorem{proposition}{Proposition}
\def \F {{\mathbb F}}
\def \P {{\mathbb P}}
\DeclareMathOperator{\ord}{ord}
\newcommand{\ra}{\rightarrow}
\title{Existence of minimal Del Pezzo surfaces of degree 1 with conic bundles over finite fields}
\author{Manoy T. Trip}
\date{Source: arXiv:2603.20054v1 (CC BY 4.0)}
\begin{document}
\maketitle

\noindent\emph{Source and licence.} Existence of minimal Del Pezzo surfaces of degree 1 with conic bundles over finite fields. Version arXiv:2603.20054v1 (CC BY 4.0), distributed by arXiv under the Creative Commons Attribution 4.0 International licence, http://creativecommons.org/licenses/by/4.0/, as recorded in the arXiv licence metadata for this exact version and shown on the abstract page https://arxiv.org/abs/2603.20054v1. Full licence notice: \texttt{licenses/arxiv-2603.20054-CC-BY-4.0.txt}.\par\smallskip

\noindent\textit{Editorial scope.} Counted unit: the article's lemma lem:type4type7 (source0.tex lines 1399--1401). The article's complete original proof of it is delivered with it (source0.tex lines 1403--1424, one \texttt{proof} environment of 22 lines). No mathematics is written, repaired or re-derived: the statement and the proof are the article's own lines, in the article's own notation, and no retained line is substituted or re-flowed.  Retained uncounted as the source-local context the proof needs: lines 857--865; lines 1031--1045.  Nothing was omitted from any retained span, so the package carries no omission record.  No retained line contains a citation, so the wrapper carries no bibliography and no bibliography entry.  No retained line contains a figure environment, an image-inclusion command or drawing code, so no image is delivered and the package declares no asset dependency.  Editorial formatting only, in addition: an AMS \texttt{amsart} preamble that restates the article's own declarations and macros used by the retained lines, namely the environments \texttt{lemma}, \texttt{proposition} on separate counters, together with the standard packages those lines need.  The retained environments print in this document as Lemma 1, Proposition 1 in order of appearance, whereas the article prints the same items with the numbering of its own full document.  The complete native source of the article as distributed in the arXiv e-print for this exact version is bundled unchanged as \texttt{sources/196824/source0.tex}.
\begin{lemma}\label{lem:galoisorbitnn/2}
    Let $\Gamma_X = \langle \tau \rangle$, where $\tau$ is the Frobenius action. Let $n = \ord(\Gamma_X)$. Then $\tau^{n/2}$ is of the form
    \begin{enumerate}
        \item $\iota_{ij}$ if $f$ is of type 1, 2 or 6,
        \item $\iota_{ijkl}$ if $f$ is of type 4 or 7,
        \item $\iota_{ijklmn}$ if $f$ is of type 3 or 5, 
    \end{enumerate}
    for distinct $i, j, k, l, m, n \in \{1, \ldots, 7\}$.
\end{lemma}

\begin{proposition}\label{prop:classificationCB}
    Let $f \colon X \ra B$ be a minimal standard rational conic bundle of degree $1$ over a perfect field $k$. Then we have one of the following cases.
    \begin{enumerate}
        \item $-K_X$ is not nef, and there is a geometrically integral bisection $A$ on $X$ such that $-K_X \cdot [A] = -1$ and $A^2 = -3$.
        \item $-K_X$ is nef but not ample, so there is an integral curve $A$ on $X$ such that $-K_X \cdot [A] = 0$. Its base change to $\overline{X}$ is of one of the following forms:
        \begin{enumerate}
            \item $A = C_1 + \cdots + C_8$, where $C_1, \ldots, C_8$ are disjoint irreducible sections on $\overline{X}$, each with self-intersection $-2$.
        
            \item $A = C_1 + \cdots + C_4$, where $C_1, \ldots, C_4$ are irreducible sections on $\overline{X}$, each with self-intersection $-2$, such that $C_i \cdot C_j = 1$ if $i \neq j$ and $i \equiv j \mod 2$, and $C_i \cdot C_j = 0$ if $i \not \equiv j \mod 2$.

            \item $A = C_1 + C_2$, where $C_1$ and $C_2$ are disjoint irreducible bisections on $\overline{X}$ with self-intersection $-2$.
        \end{enumerate}
         \item The surface $X$ is a del Pezzo surface of degree 1 admitting exactly two conic bundle structures.
    \end{enumerate}
\end{proposition}

\begin{lemma}\label{lem:type4type7}
    Let $f \colon X \ra \P^1 $ be a minimal standard conic bundle over $\F_q$ of degree 1 satisfying case (2a) in Proposition \ref{prop:classificationCB}. Then $f$ is of type 4 or 7.
\end{lemma}

\begin{proof}
    The sections $C_1, \ldots, C_8$ form a $\Gamma_X$-stable set. If $f$ is of type 3 or 5, $\Gamma_X$ contains an element of the form $\iota_{1\cdots \hat{j} \cdots 7}$ for some $j \in \{1, \ldots, 7\}$ (where $\hat{j}$ denotes the omission of the $j$-th entry) by Lemma \ref{lem:galoisorbitnn/2}. It should send $C_1$ to one of the $C_i$. We have
\begin{align*}
    \iota_{1\cdots \hat{j} \cdots 7}(C_1) = \begin{cases}
        C + F - \sum_{i=n}^7 E_n + E_1 + E_2 + E_j &\text{if} \ j \notin \{1, 2\},\\
        C + 2F - \sum_{n=1}^7 E_n + E_2 &\text{if} \ j = 1,\\
        C + 2F - \sum_{n=1}^7 E_n + E_1 &\text{if} \ j = 2.
    \end{cases}
\end{align*}
None of these possibilities is the class of one of the $C_i$, so $f$ cannot be of type 3 or 5.

If $f$ is of type $1, 2$ or $6$, $\Gamma_X$ contains an element of the form $\iota_{jk}$ for some distinct $j, k \in \{1, \ldots, 7\}$ by Lemma \ref{lem:galoisorbitnn/2}. We see that
\begin{align*}
    \iota_{jk}(C_1) = \begin{cases}
        C-F &\text{if} \ \{j, k\} = \{1, 2\},\\
        C - E_k - E_2 &\text{if} \ j = 1, k \neq 2,\\
        C - E_k - E_1 &\text{if} \ j = 2, k \neq 1,\\
        C + F - E_1 - E_2 - E_j - E_k &\text{if} \ j, k \notin \{1, 2\}.
    \end{cases}
\end{align*}
None of these are the classes of one of the $C_i$, so $f$ cannot be of type 1, 2 or 6.
\end{proof} 

\end{document}
